% Generated draft: requires mathematical and source review.
\documentclass[11pt]{article}
\usepackage[margin=25mm]{geometry}
\usepackage{fontspec}
\setmainfont{Latin Modern Roman}
\usepackage{amsmath,amssymb,mathtools,unicode-math}
\setmathfont{Latin Modern Math}
\usepackage{enumitem,needspace,xurl,hyperref}
\setlength{\parindent}{0pt}
\setlength{\parskip}{5pt}
\setlength{\emergencystretch}{3em}
\newcommand{\N}{\mathbb{N}}
\newcommand{\Z}{\mathbb{Z}}
\newcommand{\Q}{\mathbb{Q}}
\newcommand{\R}{\mathbb{R}}
\newcommand{\C}{\mathbb{C}}
\newcommand{\F}{\mathbb{F}}
\newcommand{\A}{\mathbb{A}}
\newcommand{\PP}{\mathbb{P}}
\newcommand{\E}{\mathbb{E}}
\newcommand{\eps}{\varepsilon}
\newcommand{\dd}{\,\mathrm{d}}
\begin{document}
\section*{UnsolvedMath ID: 30002438}
\textbf{Source problem code:} catalogue entry
\subsection{Definitions and mathematical statement}
For a planar set \(A\), its lattice points are the elements of \(A\cap\Z^2\). For \(R>0\) and \(c\in\R^2\), let
\[
S(c,R)=\{x\in\R^2:\|x-c\|_2=R\}.
\]
An arc \(A\) is a connected portion of this circle, and \(\operatorname{length}(A)\) denotes its Euclidean arclength.

Under the literal reading that the circle may have any center and that \(\theta\) is real, the supplied question asks whether
\[
\forall\theta\in\R,\quad \theta<1\ \Longrightarrow\ 
\exists C(\theta)\in\N\ \forall R>0\ \forall c\in\R^2\ \forall\text{ arcs }A\subseteq S(c,R),
\qquad
\operatorname{length}(A)=R^\theta
\ \Longrightarrow\ 
\#(A\cap\Z^2)\leq C(\theta).
\]
Thus the bound may depend on \(\theta\), but is uniform in the radius, the circle, and the position of the arc. The quantifier over arcs applies only where an arc of the prescribed length exists.

The center quantifier is an unresolved qualification: the supplied record does not state whether arbitrary centers, integer centers, or only the origin are intended. It also does not explicitly specify a lower bound on \(\theta\), a lower bound on \(R\), or an endpoint convention for arcs. The displayed formulation records the literal unrestricted reading; it requires confirmation against the intended source formulation.
\subsection{Short English statement}
For each fixed real exponent \(\theta<1\), is there a constant depending only on \(\theta\) that bounds the number of integer lattice points on every circle arc of length \(R^\theta\), uniformly in its radius \(R\) and position? The allowed circle centers and the intended parameter ranges require confirmation. The supplied record labels the problem partially solved; this is a statement-only draft requiring human mathematical review.
\subsection{Sources}
\begin{itemize}
\item Supplied local UnsolvedMath registry record, permanent ID \(30002438\), title ``Uniform Lattice-Point Bounds on Short Circle Arcs,'' source problem code \texttt{catalogue entry}, in the Oberwolfach Reports Open Problems collection. The record gives proposed year 2013 and supplies no named proposer.
\item Supplied source citation: \emph{Analytic Number Theory (2014)}, Oberwolfach Reports, pp.\ 2963--3037, DOI \texttt{10.4171/owr/2013/51}. \url{https://doi.org/10.4171/owr/2013/51}. The background describes these reports as documenting mathematical workshops and their open-problem discussions. This is also the literature source labeled ``original OWR report.''
\item Supplied literature source labeled ``Cilleruelo and Cordoba'': \url{https://doi.org/10.1090/S0002-9939-1992-1089403-8}.
\item Supplied literature source labeled ``Ramana conic arc bounds'': \url{https://doi.org/10.4064/aa143-3-1}.
\item Literature assessment dated exactly \texttt{2026-08-21}: the supplied curation labels the problem partially solved and states that the lattice-point assertion is proved below the square-root scale, while the full range below the radius scale remains open. This is supplied dated curation derived from the OWR material, not an independently verified assessment. The record expressly requests independent verification of time-sensitive status and cited claims.
\item The supplied statement status is \texttt{corrected\_verified}; its accompanying explanation says that mathematical content was unchanged while notation, typography, and wording were normalized into a self-contained publication form. This is a supplied editorial verification label, not independent mathematical verification here.
\end{itemize}
\end{document}
